\subsection{Generating Modules}

Let \(\mathrm{A}\) be a ring.

DEFINITION 2. --- An A-module M is called generating if every A-module N is generated by the images of the A-linear mappings from M to N.

Let \(M\) be a left A-module. We denote the dual of \(M\) by \(M^*\) and the canonical bilinear form on \(M \times M^*\) (II, §2, No.~3, p.~233) by

\[
(x, x ^ {*}) \mapsto \langle x, x ^ {*} \rangle = x ^ {*} (x).
\]

We denote by \(\tau(\mathrm{M})\) the set of elements of A of the form \(\sum_{i=1}^{n}\langle x_i,x_i^*\rangle\) , where \(x_1,\ldots ,x_n\) are elements of M and \(x_1^*,\ldots ,x_n^*\) are elements of \(\mathrm{M}^*\) . It is a two-sided ideal of A, called the trace ideal of M. The trace ideal of the A-module \(\mathrm{A}_s\) is A. The trace ideal of a module that is the direct sum of a family \((\mathrm{M}_i)_{i\in \mathrm{I}}\) of A-modules is the ideal \(\sum_{i\in \mathrm{I}}\tau(\mathrm{M}_i)\) . If M is a projective A-module, then it follows from Proposition 12 of II, §2, No.~6, p.~238, that we have \(\mathrm{M} = \tau(\mathrm{M})\mathrm{M}\) .

THEOREM 1. --- Let M be a left A-module. The following properties are equivalent:

\begin{enumerate}
\def\labelenumi{(\roman{enumi})}
\item
  The A-module M is generating.
\item
  For every A-module N, there exist a set I and a surjective A-linear mapping from \(M^{(I)}\) to N.
\item
  There exist an integer \(n \geqslant 0\) and a surjective A-linear mapping from \(\mathbf{M}^n\) to \(\mathrm{A}_s\) .
\item
  There exists an integer \(n \geqslant 0\) such that \(A_s\) is isomorphic to a direct factor submodule of \(M^n\) .
\item
  The trace ideal \(\tau (\mathbf{M})\) is equal to A.
\item
  There exist an integer \(n \geqslant 0\) , elements \(x_{1}, \ldots, x_{n}\) of M, and elements \(x_{1}^{*}, \ldots, x_{n}^{*}\) of M* satisfying \(\sum_{i=1}^{n} \langle x_{i}, x_{i}^{*} \rangle = 1\) .
\item
  \(\Rightarrow\) (ii): Suppose that M is generating, and let N be a left A-module. There exists a family \((u_{i})_{i\in\mathrm{I}}\) of A-linear mappings from M to N such that we have \(N=\sum_{i\in\mathrm{I}}u_{i}(\mathrm{M})\) . The mapping \((x_{i})\mapsto\sum_{i\in\mathrm{I}}u_{i}(x_{i})\) from \(\mathrm{M}^{(\mathrm{I})}\) to N is A-linear and surjective.
\item
  \(\Rightarrow\) (iii): Assume that property (ii) holds, and take \(N = A_s\) . We then have a set I and a surjective linear mapping \(u: M^{(I)} \to A_s\) . Since \(A_s\) is generated by the element 1, there exists a finite subset J of I such that \(u(M^{(J)}) = A_s\) , whence (iii).
\item
  \(\Rightarrow\) (iv): This follows from Proposition 21 of II, §1, No.~12, p.~218.
\item
  \(\Rightarrow\) (v): Let \(n \geqslant 1\) be an integer such that \(\mathrm{A}_s\) is isomorphic to a direct factor submodule of M. We have \(\mathrm{A} = \tau(\mathrm{A}_s) \subset \tau(\mathrm{M}^n) = \tau(\mathrm{M})\) and therefore \(\tau(\mathrm{M}) = \mathrm{A}\) .
\item
  \(\Rightarrow\) (vi): This is clear.
\item
  \(\Rightarrow\) (i): Let \(n\) be an integer \(\geqslant 0\) , \(x_{1},\ldots ,x_{n}\) elements of M, and \(x_{1}^{*},\ldots ,x_{n}^{*}\) elements of M \(^*\) satisfying \(\sum_{i = 1}^{n}\langle x_i,x_i^*\rangle = 1\) . Let N be a left A-module and \(y\) an element of N. The mappings \(u_{i}:x\mapsto \langle x,x_{i}^{*}\rangle y\) from M to N are A-linear, and we have \(y = \sum_{i = 1}^{n}u_{i}(x_{i})\) . This proves that M is a generating A-module.
\end{enumerate}

COROLLARY. --- A generating A-module is faithful.

Let \(a \in A\) satisfy \(aM = 0\) . Using the implication (i) \(\Rightarrow\) (iv) of Theorem 1, we obtain \(aA_s = 0\) , and therefore \(a = 0\) . The corollary follows (II, §1, No.~12, p.~219).

Examples. --- 1) The A-module \(A_s\) is generating.

\begin{enumerate}
\def\labelenumi{\arabic{enumi})}
\setcounter{enumi}{1}
\item
  Every nonzero free A-module is generating. More generally, every module with a generating quotient is itself generating.
\item
  Let M be a semisimple A-module whose countermodule is finitely generated. Then M is a generating \(A_{M}\) -module. Indeed, by Lemma 4 of VIII, p.~8, there exists a natural number m such that \((\mathrm{A}_{\mathrm{M}})_{s}\) is isomorphic to a submodule of \(M^{m}\) . Since \(M^{m}\) is a semisimple \(A_{M}\) -module, \((\mathrm{A}_{\mathrm{M}})_{s}\) is isomorphic to a direct factor submodule of \(M^{m}\) and M is a generating \(A_{M}\) -module (Theorem 1).
\item
  Let A be a principal ideal domain, and let P be a finitely generated A-module. There exist an integer \(n \geqslant 1\) and an increasing sequence of ideals \((\mathfrak{a}_{i})_{1 \leqslant i \leqslant n}\) of A such that P is isomorphic to the direct sum of the \(A/\mathfrak{a}_{i}\) (VII, §4, No.~4, p.~19, Theorem 2); the annihilator a of P is equal to \(a_{1}\) . Then P is a generating module over the ring A/a. If P is not a torsion module, then we have a = 0 and P is a generating A-module.
\end{enumerate}

Lemma 1. --- Let A be a commutative ring, M a finitely generated A-module, and Ann(M) its annihilator. Let \(\mathfrak{a}\) be an ideal of A. The following properties are equivalent:

\begin{enumerate}
\def\labelenumi{(\roman{enumi})}
\item
  aM = M.
\item
  Ann(M) + a = A.
\item
  There exists an element \(a\) of \(\mathfrak{a}\) such that \(am = m\) for every \(m \in \mathrm{M}\) . (i) \(\Rightarrow\) (ii): Let \((x_1, \ldots, x_n)\) be a generating family of the A-module M. If \(\mathfrak{aM} = \mathrm{M}\) , then each \(x_i\) can be written as \(\sum_{j=1}^{n} c_{ij} x_j\) , where the \(c_{ij}\) belong to \(\mathfrak{a}\) . Denote the matrix \((c_{ij})\) by \(C\) and the column matrix with entries \(x_1, \ldots, x_n\) by \(X\) . We have \((I_n - C)X = 0\) . Let \(d\) be the determinant and \(V\) the cofactor matrix of the matrix \(I_n - C\) . By formula (26) of III, §8, No.~6, p.~532, we have \(dX = {}^t V(I_n - C)X = 0\) , so that \(d \in \operatorname{Ann}(\mathrm{M})\) . On the other hand, since the \(c_{ij}\) belong to \(\mathfrak{a}\) , we have \(d \equiv 1 (\bmod \mathfrak{a})\) (III, §8, No.~5, p.~529, (18)), and therefore \(1 \in \operatorname{Ann}(\mathrm{M}) + \mathfrak{a}\) .
\item
  \(\Rightarrow\) (iii): Assuming that (ii) holds, there exist \(a\in \mathfrak{a}\) and \(b\in \mathrm{Ann}(\mathrm{M})\) such that \(a + b = 1\) . We then have \(am = m\) for every \(m\in \mathrm{M}\) .
\item
  \(\Rightarrow\) (i): This is clear.
\end{enumerate}

PROPOSITION 3. --- Let A be a commutative ring. Every finitely generated and faithful projective A-module is generating. More generally, a finitely generated projective A-module P is a generating \(A_{P}\) -module.

Let P be a finitely generated projective A-module. We have \(\tau(\mathrm{P})\mathrm{P}=\mathrm{P}\) (VIII, p.~80).

If the A-module P is faithful, then the ideal \(\tau(\mathrm{P})\) is equal to A (Lemma 1) and the A-module P is generating (Theorem 1); this leads to the first assertion. The second follows by Lemma 2 below.

Lemma 2. --- Let A be a ring and M a projective A-module. The \(\mathrm{A_M}\) -module M is projective.

Let \((x_{i})_{i\in\mathrm{I}}\) be a generating family of the A-module M. There exists a family \((x_{i}^{*})_{i\in\mathrm{I}}\) of linear forms on the A-module M such that for every \(x\in\mathrm{M}\) , the family \((\langle x,x_{i}^{*}\rangle)_{i\in\mathrm{I}}\) has finite support and \(x=\sum_{i\in\mathrm{I}}\langle x,x_{i}^{*}\rangle x_{i}\) (II, §2, No.~6, p.~238, Proposition 12). For every \(i\in\mathrm{I}\) , the mapping \(x\mapsto\langle x,x_{i}^{*}\rangle_{\mathrm{M}}\) is a linear form on the \(A_{M}\) -module M, and we have \(x=\sum_{i\in\mathrm{I}}\langle x,x_{i}^{*}\rangle_{\mathrm{M}}x_{i}\) for every \(x\in\mathrm{M}\) . By loc. cit., M is a projective \(A_{M}\) -module.

